\section{The model: stretch probability rank}
\label{sec:model}

\subsection{Quantiles and transport speed}

Let $Q(u)$ be the quantile function of $X$.  If $U$ is uniform on $(0,1)$,
then $Q(U)$ has that law.  Write $q(u)=Q'(u)$ for the quantile density.
Differentiating $F_X(Q(u))=u$ gives
\begin{equation}
f_X(Q(u))=\frac1{q(u)}.
\label{eq:reciprocal}
\end{equation}
The derivative $q$ measures log-return distance per unit of percentile rank.
Small $q$ compresses many ranks into a narrow return interval and creates
high density; large $q$ spreads the same probability over a wider interval
and creates low density.  A strictly positive $q$ makes $Q$ strictly
increasing.  Parameterizing $\log q$ makes that positivity automatic.

The original LQD model starts from
\begin{equation}
q(u)=\frac{e^{g(u)}}{u(1-u)}.
\label{eq:originalskeleton}
\end{equation}
The factors $u(1-u)$ make $q$ diverge at both endpoints, giving unbounded
support.  They also cancel the derivative of the logistic map.  Since
$\dd u/\dd z=u(1-u)$,
\begin{equation}
\frac{\dd Q}{\dd z}=e^{g(\Lambda(z))}.
\label{eq:originalspeed}
\end{equation}
Thus $g$ is the logarithm of the speed at which log-return moves along the
log-odds line.  If $g$ tends to a constant in a tail, that speed tends to a
positive constant and the log-return tail is exponential.

To permit lighter tails, slow the transport gradually at extreme ranks.
Introduce the positive logarithmic gauges
\begin{equation}
\ell_-(u)=1-\log u,
\qquad
\ell_+(u)=1-\log(1-u).
\label{eq:ellgauges}
\end{equation}
They behave like $\log(1/u)$ on the left and
$\log(1/(1-u))$ on the right.  Choose tail exponents
$\alpha_-,\alpha_+\in[0,1/2]$ and set
\begin{equation}
\log q(u)=
-\log u-\log(1-u)
-\alpha_-\log\ell_-(u)-\alpha_+\log\ell_+(u)+g(u).
\label{eq:skeleton}
\end{equation}
Equivalently,
\begin{equation}
q(u)=\frac{e^{g(u)}}
{u(1-u)\ell_-(u)^{\alpha_-}\ell_+(u)^{\alpha_+}}.
\label{eq:lqdform}
\end{equation}
This replaces the left endpoint factor $u$ by
$u\{\log(1/u)\}^{\alpha_-}$ to leading order, and does the analogous thing
on the right.

Because $\dd u/\dd z=u(1-u)$, the two leading endpoint singularities cancel:
\begin{equation}
\frac{\dd Q}{\dd z}
=\frac{e^{g(\Lambda(z))}}
{\ell_-(\Lambda(z))^{\alpha_-}
 \ell_+(\Lambda(z))^{\alpha_+}}.
\label{eq:speed}
\end{equation}
The endpoint singularities have disappeared.  The distribution is obtained
by integrating an ordinary positive function on the real line.  Setting
$\alpha_-\!=\!\alpha_+\!=\!0$ recovers the original model.  Increasing one
exponent slows the corresponding extreme transport and lightens that tail.

\begin{figure}[htbp]
\centering
\paperfig{\textwidth}{fig_transport.pdf}
\caption{The transport construction in the exponential-tail subclass with
constant speed.  Left: logistic log-odds and five percentile marks.  Center:
an increasing map from log-odds to log-return; strike roots are read on this
map.  Right: the implied-volatility smile obtained by pricing the resulting
law.  The first two panels define the model; the last changes coordinates.}
\label{fig:transport}
\end{figure}

\subsection{A complete slice}

The function $g$ controls the body of the law.  A low-degree polynomial is
enough for smooth smiles and keeps the number of effective shapes visible.
Shifted Legendre polynomials are orthogonal under uniform rank, so their
coefficients are substantially better scaled than raw powers of $u$.

\begin{definition}[Generalized LQD slice]
\label{def:lqd}
Fix $N\ge2$, tail exponents $\alpha_\pm\in[0,1/2]$, and coefficients
$\theta=(L,R,a_2,\ldots,a_N)$.  Define
\begin{equation}
g(u)=(1-u)L+uR+\sum_{n=2}^N a_nP_n(1-2u).
\label{eq:gpoly}
\end{equation}
\begin{equation}
x'(z)=\frac{e^{g(\Lambda(z))}}
{\ell_-(\Lambda(z))^{\alpha_-}
 \ell_+(\Lambda(z))^{\alpha_+}}.
\label{eq:xspeed}
\end{equation}
\begin{equation}
\boxed{\quad X=x(Z),\qquad
x(z)=m+\int_0^z x'(t)\,\dd t,\qquad Z\sim\Lambda.\quad}
\label{eq:lqd}
\end{equation}
The shift $m$ is chosen, when possible, so that $\E[e^X]=1$.
\end{definition}

For every finite coefficient vector, $x'(z)>0$.  Since $\alpha_\pm<1$, the
integral of $x'$ diverges at both ends, so $x$ is a smooth increasing
bijection from $\R$ to $\R$.  Consequently
\begin{equation}
F_X(x(z))=\Lambda(z),\qquad
f_X(x(z))=\frac{\rho(z)}{x'(z)}>0.
\label{eq:density}
\end{equation}
The map cannot fold, stop, or reverse.  The density is positive before any
market quote is fitted.  The parameters now have three distinct roles:
\begin{itemize}
\item $\alpha_\pm$ select the two asymptotic tail classes;
\item the endpoint values of $g$ set the two tail scales;
\item the remaining coefficients shape the distribution between the tails.
\end{itemize}

The endpoint coefficients are
\begin{equation}
\lambda_-=e^{g(0)}
=\exp\left(L+\sum_{n\ge2}a_n\right),\qquad
\lambda_+=e^{g(1)}
=\exp\left(R+\sum_{n\ge2}(-1)^na_n\right).
\label{eq:endpoints}
\end{equation}
When $\alpha_\pm=0$, these numbers are the limiting transport speeds.  For
$\alpha_\pm>0$, the speeds tend to zero and $\lambda_\pm$ multiply their
power-law decay in $|z|$.  In both cases $\lambda_\pm$ are tail scales, while
$\alpha_\pm$ determine the type of tail.

\begin{figure}[htbp]
\centering
\paperfig{\textwidth}{fig_modes.pdf}
\caption{Three Legendre perturbations of one exponential-tail reference law.
Left: their action on the log-speed $g$.  Center: the implied-volatility
response.  Right: the density response.  Raising $g$ stretches return space
and lowers density at the affected ranks.  Even modes act on both sides; odd
modes transfer shape from one side to the other.}
\label{fig:modes}
\end{figure}

Raw Legendre coefficients move both the body and the endpoint scales.  That
mixing is inconvenient in calibration.  Replace them by the
endpoint-vanishing modes
\begin{equation}
\widetilde P_n(u)=P_n(1-2u)-(1-u)-(-1)^nu,\qquad n\ge2.
\label{eq:endpointbasis}
\end{equation}
They satisfy $\widetilde P_n(0)=\widetilde P_n(1)=0$.  The equivalent
parameter vector $(\log\lambda_-,\log\lambda_+,a_2,\ldots,a_N)$ lets an
optimizer change body shape without silently changing the asymptotic tails.

\subsection{One shift matches the forward}

The transport fixes the shape of the log-return law but leaves its location
free.  Translating every log-return by the same number changes neither ranks,
density shape, nor tail class.  Use this one degree of freedom to impose
$\E[e^X]=1$.

Let
\[
\bar x(z)=\int_0^zx'(t)\,\dd t,
\qquad
I=\int_{-\infty}^{\infty}e^{\bar x(z)}\rho(z)\,\dd z.
\]

\begin{proposition}[The forward condition]
\label{prop:shift}
If $I<\infty$, the unique martingale shift is $m=-\log I$.  Moreover,
\begin{equation}
I<\infty
\quad\Longleftrightarrow\quad
\alpha_+>0\ \text{or}\ (\alpha_+=0\ \text{and}\ \lambda_+<1).
\label{eq:martcondition}
\end{equation}
\end{proposition}

\begin{proof}
Since $X=m+\bar x(Z)$, $\E[e^X]=e^mI$, which gives the unique shift.
If $\alpha_+>0$, then $\bar x(z)=O(z^{1-\alpha_+})=o(z)$ while
$\rho(z)\sim e^{-z}$, so the right integrand is integrable.  If
$\alpha_+=0$, then $\bar x(z)=\lambda_+z+O(1)$ and the integrand is
asymptotic to a constant times $e^{(\lambda_+-1)z}$.  The left integrand is
always integrable because $\bar x(z)\le0$ for $z\le0$ and
$\rho(z)\sim e^z$.
\end{proof}

Only the right tail can make the forward infinite.  In the exponential-tail
subclass its scale must satisfy the exact wall $\lambda_+<1$.  Every strictly
lighter right tail has a finite forward for any finite $\lambda_+$.  The jump
at $\alpha_+=0$ is mathematical, not conventional: any positive exponent
eventually beats a linear exponential rate.  The convergence is slow when
the exponent is small, which matters numerically in \cref{sec:numerics}.

\subsection{Two reference laws anchor the model}
\label{sec:reference-laws}

Set $\alpha_\pm=0$, $L=R=\log s$, and $a_n=0$, with $0<s<1$.  Then
$x(z)=m+sz$ and $X$ is logistic with scale $s$.  Its moment-generating
function is
\begin{equation}
\E[e^{tZ}]
=\int_0^1\left(\frac{u}{1-u}\right)^t\dd u
=\Gamma(1+t)\Gamma(1-t)
=\frac{\pi t}{\sin(\pi t)},\qquad |t|<1.
\label{eq:mgf}
\end{equation}
Hence
\begin{equation}
m=-\log\frac{\pi s}{\sin(\pi s)},
\qquad \operatorname{Var}(X)=\frac{\pi^2s^2}{3}.
\label{eq:logisticreference}
\end{equation}
Every quantity is explicit, so this law is the primary end-to-end test of an
implementation: transport, normalization, moments, strike roots, prices, and
densities must all agree with it.

At the lighter endpoint of the tail range, a normal law
$X\sim N(-s^2/2,s^2)$ has
\[
q(u)\sim\frac{s}{u\sqrt{2\log(1/u)}}\quad(u\downarrow0),
\]
and the symmetric right relation.  With $\alpha_\pm=1/2$, the function $g$
obtained from \eqref{eq:skeleton} has finite endpoint limits.  Polynomial
LQD log-speeds can therefore approximate the normal law without forcing
$g$ to diverge at either endpoint.  The logistic and normal laws are not
competing calibrations; they are exact benchmarks for the two ends of the
tail range developed next.

\begin{figure}[htbp]
\centering
\paperfig{0.92\textwidth}{fig_exact.pdf}
\caption{The constant-speed logistic slice as an implementation test.  Left:
computed transport and martingale shift errors against
\eqref{eq:logisticreference}.  Right: the variance-matched cold start and its
small-variance ATM bias.  The test exercises transport integration,
normalization, tail treatment, and strike inversion with known answers.}
\label{fig:exact}
\end{figure}
